\section{Fraud Proof Basics}


The fraud proof protocol enables efficient verification of disputed state transitions through a carefully designed challenge-response mechanism.

\subsection{Assertion Protocol}

When a sequencer makes an assertion about a state transition, they essentially claim:

\begin{proposition}[State Transition Assertion]
Given initial state $S_0$ and transaction sequence $\mathcal{T} = \{t_1, t_2, ..., t_n\}$, applying the state transition function yields final state $S_n$:
\begin{equation}
S_n = \text{STF}^n(S_0, \mathcal{T})
\end{equation}
where $\text{STF}^n$ denotes $n$ sequential applications of the state transition function.
\end{proposition}

\subsection{Challenge Mechanism}

Any party observing an invalid assertion can initiate a challenge by:

\begin{enumerate}
    \item Computing the correct final state $S'_n$
    \item Demonstrating $S'_n \neq S_n$
    \item Staking collateral to back the challenge
    \item Initiating the interactive dispute resolution protocol
\end{enumerate}

\subsection{Interactive Proof Protocol}

The proof protocol follows an interactive game where challenger and defender exchange messages to narrow the dispute:

\begin{algorithm}
\caption{Interactive Fraud Proof Protocol}
\begin{algorithmic}[1]
\State \textbf{Input:} Assertion $A = (S_0, S_n, \mathcal{T})$, Challenger $C$, Defender $D$
\State \textbf{Output:} Winner $\in \{C, D\}$
\State $\text{left} \gets 0$, $\text{right} \gets n$
\While{$\text{right} - \text{left} > 1$}
    \State $\text{mid} \gets \lfloor(\text{left} + \text{right})/2\rfloor$
    \State $D$ provides $S_{\text{mid}}$
    \State $C$ computes expected $S'_{\text{mid}}$
    \If{$S_{\text{mid}} = S'_{\text{mid}}$}
        \State $\text{left} \gets \text{mid}$ \Comment{Disagreement in right half}
    \Else
        \State $\text{right} \gets \text{mid}$ \Comment{Disagreement in left half}
    \EndIf
\EndWhile
\State Execute single step $\text{left} \rightarrow \text{right}$ on-chain
\State \textbf{return} Winner based on execution result
\end{algorithmic}
\end{algorithm}

\subsection{On-chain Resolution}

Once the dispute is narrowed to a single step, the Layer 1 contract executes that step to determine the winner:

\begin{equation}
\text{Winner} = \begin{cases}
\text{Defender} & \text{if } S_{\text{right}} = \text{STF}(S_{\text{left}}, t_{\text{left}}) \\
\text{Challenger} & \text{otherwise}
\end{cases}
\end{equation}

